\documentclass[11pt,a4paper]{article}
\usepackage[T1]{fontenc}
\usepackage{lmodern,amsmath,amssymb,amsthm,mathtools,microtype}
\usepackage[margin=27mm]{geometry}
\usepackage[colorlinks=true,linkcolor=blue,citecolor=blue,urlcolor=blue]{hyperref}
\setlength{\emergencystretch}{2em}
\newtheorem{theorem}{Theorem}[section]
\newtheorem{lemma}[theorem]{Lemma}
\newcommand{\Q}{\mathbb Q}
\DeclareMathOperator{\Tr}{Tr}
\title{An intrinsic congruent-curve pole block in the general elliptic source}
\author{}\date{17 November 2026}
\begin{document}\maketitle
\begin{abstract}
We transfer the intrinsic two-coordinate congruent-curve pole equation
to the normalized general elliptic source. A coherent change of the
physical weight gives its required real range and two-adic class,
while preserving the local extraction and height proof. The new pole
equation has no arithmetic nonzero guards. Its terminal radicand may
be made a unit at every point-algebra character by the proved density
of its actual rational witnesses. The transfer uses Burungale's
simultaneous-twist theorem in place of the former pole realization
premise, retains all other stated source premises, and has twenty-five
total variables. No numerical degree bound or smaller variable count
is asserted.
\end{abstract}

\section{Specification and premises}

Use either the normalized general source of \cite{Reweighted} or its
order-two scalar revision \cite{Precision}. Retain its nine point
roles, integer-reflecting input label, scalar \(\rho=\rho_L(X)\),
common slope, two contact functions \(R_j\), unsquared comparison
scales, tests and assigned orders. Either the original scalar or
the aligned scalar of \cite{Aligned} may be used.
The fixed exceptional set contains \(2,7\).

Every integer input has the following established scalar properties:
\[
 \rho>0,\qquad \rho\text{ has square reduced denominator},\qquad
 \rho\in\Q_2^{\times2},\qquad v_2(\rho)\ge4.
\]
The last valuation is four for the original scalar and twelve for
the aligned scalar. These statements are used only in completeness.
The group, prime-pattern, normalized contact-height and aggregate
trace-height premises retain exactly their scope in the cited sources.
The old pole realization premise is replaced by Burungale's Theorem B
in the precise form proved sufficient in \cite{Intrinsic}.

In both blocks put
\begin{equation}\label{eq:weights}
 u_j=\rho R_j,\qquad w_j=\frac{u_j}{(1+7u_j)^2}.
\end{equation}
Use this physical weight in \emph{every} field and scalar detector
denominator. Thus, with the unchanged normalized scalar numerators,
\[
 W_j=\prod_i\left(1+\frac{C_i^2}
                           {\delta_j^2w_j^{2a_i-1}}\right)
     \prod_k\left(1+\frac{G_{j,k}^2}
                           {w_j^{2\kappa_k-1}}\right),
 \qquad
 G_{j,k}=R_j^{\kappa_k+1}\frac{R_jS_k}{1+R_jS_k}.
\]
Retain all point, difference, scalar and trace denominators, and
guard \(1+7u_j\). Put
\[
 V_j=\Tr(W_j^{-1}),\qquad b_j=\frac{4w_j}{V_j^2},\qquad d_j=7b_j.
\]
Replace each old pole pair by two unrestricted rational coordinates
\(v_j,s_j\) and the single finite equation
\begin{equation}\label{eq:pole}
 Y_j^2=P_{d_j}(v_j,s_j),\qquad
 P_d(v,s)=d(d+v^2)^2s^4-d(d-v^2)^2(1-d)^2.
\end{equation}
No arithmetic guard on \(s_j,v_j,Y_j,d_j\pm v_j^2,1-d_j\) is imposed
by this pole equation. For the separable one-coordinate encoding we
retain the unit guard on the cleared terminal radicand. The density
argument below supplies it without a further free coordinate.

\section{Local soundness}

\begin{lemma}\label{lem:weight}
For \(p\ne7\), every odd valuation of \(w_j\) is positive.
More precisely, if \(n=v_p(u_j)\) is odd, then
\(v_p(w_j)=|n|\). For positive real \(u_j\), one has
\(0<w_j\le1/28\).
\end{lemma}
\begin{proof}
For nonzero \(n\), the denominator in \eqref{eq:weights} has
valuation \(2\min(0,n)\). If \(n=0\), the weight valuation is
even. Finally \((1+7u)^2-28u=(1-7u)^2\ge0\).
\end{proof}

\begin{lemma}\label{lem:pole}
Equation \eqref{eq:pole} excludes negative odd valuations of \(d\)
at every prime, without any nonzero hypotheses.
\end{lemma}
\begin{proof}
Suppose \(k=v_p(d)<0\) is odd and set
\(m=\min(k,2v_p(v))\). The two summands of \(d\pm v^2\)
have different valuations, so both have valuation \(m\), and
\(v_p(1-d)=k\). For \(s\ne0\), the two radicand terms have
valuations
\[
 k+2m+4v_p(s),\qquad 3k+2m.
\]
Both are odd and they are distinct, because \(2v_p(s)\ne k\).
Their sum therefore has odd valuation. For \(s=0\) the second
term alone has odd valuation. The same calculation includes \(v=0\)
by taking its valuation to be infinity. A rational square cannot
have the resulting odd valuation, including when \(Y=0\).
\end{proof}

At every prime outside the unchanged exceptional set, Lemma~\ref{lem:pole}
for \(d_j=7b_j\) excludes the same negative odd valuations of \(b_j\)
as the former pole predicate. Lemma~\ref{lem:weight} gives exactly
the normalized weight parity used in \cite{Reweighted,Precision}.
The trace-factor extraction, scalar divisibility, comparison
divisibility and height proof consequently apply unchanged.
This includes the order-two revision's relative logarithmic precision
eight, sixth-order comparisons and squared Lucas parameter.
Thus its soundness conclusion holds for every sufficiently large
fixed odd \(L\), with the same scope of the remaining premises.

\section{Complete preparation of the selected rational inputs}

Fix an integer input and a fixed odd \(L\ge3\).
After fixing \(\rho\), first form the temporary finite preparation
and avoidance set, including the prime supports of \(\rho\) and
\(1+7\rho\) and all inherited fixed data. This is a completeness
choice, not an enlargement of the global exceptional set used in
the height proof. Choose the reference, scalar and root points in
the retained order using this set, then fix all comparison scales.
The scalar functions can be made arbitrarily
small in the required real and fixed finite neighborhoods.
The limiting physical weights are now
\[
                 w_{j,0}=\frac{\rho}{(1+7\rho)^2}\ne0.
\]
At each fixed place, \(R_j\to1\), the field attenuation still tends
to zero, and the unchanged scalar preparation makes the actual
detectors as close to one as needed. Keep every inherited fixed
avoidance and denominator condition, enlarging only the temporary
completeness set when the later fixed root data require it.

At selected fresh contacts, \(p\ne7\), \(v_p(u_j)=1\), and
\(1+7u_j\) is a unit. All field and scalar local estimates and
the original primary extra divisibility are therefore unchanged.
The reciprocal trace is a unit and \(v_p(b_j)=1\).
The original contact factorization shows that every other odd
valuation of the raw weight is at a fixed preparation prime.
The new denominator \((1+7u_j)^2\) and the trace square change
only even valuations. Preserving the prepared square classes at
the fixed primes therefore gives square reduced denominators
for both \(b_j\).

For the actual real character, choose both slots of \(W_j\) in
\([1,2)\). Then \(V_j>1\), and Lemma~\ref{lem:weight} gives
\[
                    0<b_j=4w_j/V_j^2<1/7.
\]
Only the selected rational character needs this bound.

At two, the sufficiently close contact and detector preparation gives
\[
 b_j\in\Q_2^{\times2},\qquad
 v_2(b_j)=v_2(\rho)\ge4.
\]
Indeed \(R_j\) is a local square close to one,
\(1+7u_j\) is a unit, and \(V_j\) is sufficiently close to two
to have valuation one. The quotient \(b_j/u_j\) is a rational square.

At seven let \(n=v_7(\rho)\). Its square reduced denominator
implies that either \(n\ge0\) or \(n\le-2\) is even. Directly,
\[
 v_7\left(\frac{\rho}{(1+7\rho)^2}\right)
 =\begin{cases}n,&n\ge0,\\-n-2,&n\le-2\text{ even}.
 \end{cases}
\]
Both alternatives are nonnegative. The sufficiently close preparation,
with \(V_j\) a seven-adic unit, gives \(v_7(b_j)\ge0\).
Consequently multiplication by seven preserves the square reduced
denominator, and
\[
 0<d_j<1,\qquad d_j\in7\Q_2^{\times2},\qquad v_2(d_j)\ge4.
\]
These are exactly the hypotheses of the intrinsic density theorem
in \cite{Intrinsic}. It supplies rational witnesses to each equation
\eqref{eq:pole}.

\section{Every character and the coordinate count}

The retained source preparation makes all point-algebra characters
real. At every such character the physical weights are positive,
all detector factors are at least one, and the reciprocal traces
are positive. Thus \(d_j\) is a unit of the point algebra, although
its values at other characters need not lie below one.
This causes no restriction: equation \eqref{eq:pole} has no
\(1-d_j\) guard and its local implication already includes \(d_j=1\).

At a geometric character taking \(d_j\) to \(c\ne0\), the
polynomial \(P_c(v,s)\) is nonzero, since its leading coefficient
in \(s\) is \(c(c+v^2)^2\). Its norm from the point algebra is
therefore a nonzero polynomial in \(v,s\). The intrinsic theorem
gives a Zariski-dense set of actual rational witness pairs, so one
can avoid this norm and the finitely many retained denominator
norms simultaneously. The cleared radicand is consequently a
whole-algebra unit. Each terminal quadratic extension is etale,
and the existing one-coordinate unit-norm encoder applies directly.

\begin{theorem}
Under the retained group, normalized height, trace-height and
prime-pattern premises and Burungale's Theorem B in the stated
intrinsic form, the modified general source defines the integers
for every sufficiently large fixed odd \(L\). It is complete
for every integer input and every fixed odd \(L\ge3\).
Its polynomial specification has twenty-five total variables:
the input, eighteen point coordinates, the common slope, four
pole coordinates, and the one encoder coordinate.
\end{theorem}
\begin{proof}
Soundness follows from the unchanged extraction and height arguments
and Lemmas~\ref{lem:weight}--\ref{lem:pole}. Ordered preparation,
the intrinsic realization theorem, and simultaneous guard avoidance
prove completeness. All added quantities besides the four displayed
pole coordinates are rational source functions or finite monic
generators, so the stated census is exact.
\end{proof}

No numerical sufficient value of \(L\), degree bound for the complete
determinant, or smaller defining-variable bound is asserted here.

\begin{thebibliography}{9}\raggedright
\bibitem{Reweighted}\emph{Quadratic reweightings of native elliptic detectors},
\nolinkurl{papers/pole_parameters/reweighted_elliptic_detectors_20261116_v53.tex}.
\bibitem{Precision}\emph{Scalar precision two with an eighth-order elliptic comparison},
\nolinkurl{precision_transfer_20261117_v56.tex}.
\bibitem{Aligned}\emph{A scalar norm from the detector field},
\nolinkurl{papers/pole_parameters/detector_field_norm_scalar_20261116_v52.tex}.
\bibitem{Intrinsic}\emph{A two-coordinate pole test of coefficient height four},
\nolinkurl{intrinsic_burungale_pole.tex}.
\bibitem{Burungale}A.~Burungale,
\emph{A Diophantine definition of $\mathbb Z$ in $\mathbb Q$},
arXiv:2611.16130v1, 12 November 2026, Theorem B.
\url{https://arxiv.org/abs/2611.16130}.
\end{thebibliography}
\end{document}
